
tôt de celle des carrés correspondants (7.10.0), compte tenu de la transitivité des images inverses (IV 2.12).

\textbf{Remarque 7.10.3.} Lorsque @@M0@@ est propre, la commutativité de (7.10.1) implique aussitôt celle du diagramme suivant, qui s'en déduit par passage aux gradués associés :
@@M1@@
De même, si @@M2@@ et @@M3@@ sont à fibres régulières et séparés, la commutativité de (7.10.2) implique celle de
@@M4@@
Lorsque @@M5@@ et @@M6@@ sont de plus quasi-projectifs et lisses sur @@M7@@, on sait d'ailleurs que @@M8@@ et @@M9@@ sont compatibles aux filtrations par la codimension des supports, et on trouve un diagramme commutatif comme (7.10.2 bis), mais les @@M10@@ étant remplacés par les @@M11@@ (cf. Exp. 0, 2.1). On peut aussi noter que si @@M12@@ et @@M13@@ sont quasi-projectifs sur @@M14@@ à fibres régulières et séparés, alors @@M15@@ et @@M16@@ sont automatiquement des morphismes d'intersection complète (VIII 1.1), et on obtient un diagramme commutatif analogue à (7.10.1 bis), où les @@M17@@ sont remplacés par les @@M18@@, compte tenu du résultat de compatibilité (VIII 3.2).

